\newpage
\phantomsection
\label{prf:finite-linear-combination}

\begin{remark*}[Return]
\hyperref[thm:finite-linear-combination-rule]{Return to Theorem}
\end{remark*}

\begin{theorem*}[Finite Linear Combination Rule]
Let $I \subseteq \mathbb{R}$ be an open interval, let $c \in I$, let $n \in \mathbb{N}$, let $f_1, \ldots, f_n : I \to \mathbb{R}$ be differentiable at $c$, and let $\alpha_1, \ldots, \alpha_n \in \mathbb{R}$. Then $\sum_{i=1}^{n} \alpha_i f_i$ is differentiable at $c$ and
\[
\left(\sum_{i=1}^{n} \alpha_i f_i\right)'(c) = \sum_{i=1}^{n} \alpha_i f_i'(c).
\]
\end{theorem*}

\begin{proof}
\textbf{Professional Standard Proof.}~\\
The proof proceeds by induction on $n$. For $n = 1$, the result follows from the Constant Multiple Rule. For the inductive step, assume the statement holds for $n = k$ and consider $n = k+1$. Write $\sum_{i=1}^{k+1} \alpha_i f_i = \left(\sum_{i=1}^{k} \alpha_i f_i\right) + \alpha_{k+1} f_{k+1}$. By the inductive hypothesis, $\sum_{i=1}^{k} \alpha_i f_i$ is differentiable at $c$ with derivative $\sum_{i=1}^{k} \alpha_i f_i'(c)$. By the Constant Multiple Rule, $\alpha_{k+1} f_{k+1}$ is differentiable at $c$ with derivative $\alpha_{k+1} f_{k+1}'(c)$. The Sum Rule then gives the desired result.
\end{proof}

\begin{proof}
\textbf{Detailed Learning Proof.}~\\

\textbf{Step 1.} Base case $n = 1$. Define $h = \alpha_1 f_1$. By the Constant Multiple Rule, $h$ is differentiable at $c$ and $h'(c) = \alpha_1 f_1'(c)$. This equals $\sum_{i=1}^{1} \alpha_i f_i'(c)$, establishing the base case.

\textbf{Step 2.} Inductive hypothesis. Assume the statement holds for some $n = k \geq 1$: for any $k$ functions differentiable at $c$ and any scalars $\alpha_1, \ldots, \alpha_k \in \mathbb{R}$,
\[
\left(\sum_{i=1}^{k} \alpha_i f_i\right)'(c) = \sum_{i=1}^{k} \alpha_i f_i'(c).
\]

\textbf{Step 3.} Inductive step for $n = k+1$. Consider $h = \sum_{i=1}^{k+1} \alpha_i f_i$. Write this as $h = \left(\sum_{i=1}^{k} \alpha_i f_i\right) + \alpha_{k+1} f_{k+1}$. Let $g = \sum_{i=1}^{k} \alpha_i f_i$. By the inductive hypothesis, $g$ is differentiable at $c$ with $g'(c) = \sum_{i=1}^{k} \alpha_i f_i'(c)$. By the Constant Multiple Rule, $\alpha_{k+1} f_{k+1}$ is differentiable at $c$ with derivative $\alpha_{k+1} f_{k+1}'(c)$.

\textbf{Step 4.} Application of the Sum Rule. Since $h = g + \alpha_{k+1} f_{k+1}$ where both $g$ and $\alpha_{k+1} f_{k+1}$ are differentiable at $c$, the Sum Rule gives
\[
h'(c) = g'(c) + \alpha_{k+1} f_{k+1}'(c) = \sum_{i=1}^{k} \alpha_i f_i'(c) + \alpha_{k+1} f_{k+1}'(c) = \sum_{i=1}^{k+1} \alpha_i f_i'(c).
\]

\textbf{Step 5.} Conclusion by induction. The statement holds for all $n \in \mathbb{N}$.
\end{proof}

\begin{remark*}[Proof structure]
The proof uses induction on $n$, establishing linearity of differentiation over finite linear combinations. The base case invokes the Constant Multiple Rule, while the inductive step decomposes the $(k+1)$-term sum by peeling off the last term, applying the inductive hypothesis to the first $k$ terms, the Constant Multiple Rule to the final term, and the Sum Rule to combine them. This structure reveals that finite linearity builds systematically from the two fundamental linearity properties: scalar multiplication and addition.
\end{remark*}

\begin{remark*}[Dependencies]~\\
\begin{itemize}
  \item \hyperref[thm:constant-multiple-rule]{Constant Multiple Rule}
  \item \hyperref[thm:sum-rule]{Sum Rule}
\end{itemize}
\end{remark*}